\section*{科研经历}
以下介绍我已经建立和正在探索的研究方向，并为每个方向列出具有代表性的论文、项目和报告。
\listsnote
\researchplannote

\subsection*{轨迹表示学习}
轨迹表示学习是我已经建立的研究方向之一，旨在学习能够通用于多种下游任务的车辆、人员或船舶轨迹表示。
该方向也包括构建可同时处理多种任务的轨迹多任务模型或基础模型。

\pubentry{IEEE TKDE, 2025, 影响因子 11.6}{\href{https://ieeexplore.ieee.org/document/10818577}{论文} \,/\, \href{https://github.com/Logan-Lin/UniTE}{代码}}
{UniTE: A Survey and Unified Pipeline for Pre-training Spatiotemporal Trajectory Embeddings}
{\me, Zeyu Zhou, Yicheng Liu, Haochen Lv, Haomin Wen, Tianyi Li, Yushuai Li, Christian S. Jensen, Shengnan Guo, Youfang Lin, Huaiyu Wan}
{综述可复用移动轨迹表示的学习方法，并用共享代码将这些方法整合为统一的模块化框架，使新方法可以在同一基础上构建、比较和评估。}

\pubentry{IEEE TKDE, 2025, 影响因子 11.6}{\href{https://ieeexplore.ieee.org/document/11004614}{论文} \,/\, \href{https://github.com/Logan-Lin/UVTM}{代码}}
{UVTM: Universal Vehicle Trajectory Modeling with ST Feature Domain Generation}
{\me, Jilin Hu, Shengnan Guo, Bin Yang, Christian S. Jensen, Youfang Lin, Huaiyu Wan}
{用单一模型处理出行时间估计、轨迹恢复和轨迹预测等多种车辆 GPS 轨迹任务，无须为每项任务维护独立模型。该方法通过将不完整轨迹重建为稠密且完整的轨迹，在轨迹稀疏或仅有部分特征时仍能保持较高精度。}

\pubentry{NeurIPS, 2025, CCF A 类会议}{\href{https://neurips.cc/virtual/2025/oral/117509}{口头报告} \,/\, \href{https://neurips.cc/virtual/2025/poster/117508}{海报}}
{TransferTraj: A Vehicle Trajectory Learning Model for Region and Task Transferability}
{Tonglong Wei*, \me*, Zeyu Zhou, Haomin Wen, Jilin Hu, Shengnan Guo, Youfang Lin, Gao Cong, Huaiyu Wan}
{从车辆 GPS 轨迹中学习可在不同区域和不同预测任务之间迁移的知识，无须重新训练或维护多个专用模型。该方法将各项任务统一转化为恢复轨迹缺失部分的问题，使单一模型在数据有限时也能处理多种任务。}

\pubentry{IEEE TKDE, 2023, 影响因子 8.9}{\href{https://ieeexplore.ieee.org/abstract/document/10375102}{论文} \,/\, \href{https://github.com/Logan-Lin/MMTEC}{代码}}
{Pre-training General Trajectory Embeddings with Maximum Multi-view Entropy Coding}
{\me, Huaiyu Wan, Shengnan Guo, Jilin Hu, Christian S. Jensen, Youfang Lin}
{从无标注数据中学习通用移动轨迹表示，同时捕捉出行行为以及空间和时间模式。所得表示避免针对特定任务的偏差，因而能够迁移到多种下游任务。}

\pubentry{AAAI, 2021, CCF A 类会议}{\href{https://ojs.aaai.org/index.php/AAAI/article/view/16548}{论文} \,/\, \href{https://github.com/Logan-Lin/CTLE}{代码}}
{Pre-training Context and Time Aware Location Embeddings from Spatial-Temporal Trajectories for User Next Location Prediction}
{\me, Huaiyu Wan, Shengnan Guo, Youfang Lin}
{从移动轨迹中预训练地点表示，捕捉地点含义随周边环境和访问时间产生的变化，从而更准确地预测用户的下一个地点。}

\subsection*{时空数据挖掘}
时空数据挖掘是包含上述研究的更广泛方向，研究如何从带有空间信息且随时间变化的大规模数据中提取并利用规律。
这类数据通常来自交通场景。

\pubentry{IJCAI, 2026, CCF B 类会议}{\href{https://arxiv.org/abs/2603.29837}{预印本} \,/\, \href{https://github.com/Logan-Lin/DiSGMM-code}{代码}}
{DiSGMM: A Method for Time-varying Microscopic Weight Completion on Road Networks}
{\me, Jilin Hu, Shengnan Guo, Christian S. Jensen, Youfang Lin, Huaiyu Wan}
{在路段之间和单个路段内部的观测均较稀疏时，补全路网中细粒度且随时间变化的交通状态，如特定时段各路段的行驶速度。该方法估计每个路段和时段可能出现的交通状态的完整范围，而非单一数值。}

\pubentry{NeurIPS, 2025, CCF A 类会议}{\href{https://neurips.cc/virtual/2025/poster/117945}{海报} \,/\, \href{https://arxiv.org/abs/2410.14281}{预印本}}
{PLMTrajRec: A Scalable and Generalizable Trajectory Recovery Method with Pre-trained Language Models}
{Tonglong Wei*, \me*, Youfang Lin, Shengnan Guo, Jilin Hu, Haitao Yuan, Gao Cong, Huaiyu Wan}
{恢复稀疏移动轨迹中的缺失点以重建详细路径，仅需少量稠密训练数据，并能泛化到不同采样率记录的轨迹。}

\pubentry{WWW, 2025, CCF A 类会议}{\href{https://openreview.net/forum?id=KmMSQS6tFn}{论文} \,/\, \href{https://github.com/decisionintelligence/Path-LLM}{代码}}
{Path-LLM: A Multi-Modal Path Representation Learning by Aligning and Fusing with Large Language Models}
{Yongfu Wei*, \me*, Hongfan Gao, Ronghui Xu, Sean Bin Yang, Jilin Hu}
{结合路网结构与描述物理环境和区域背景的文本，学习路网中的路径表示，补充了以往方法忽略的上下文信息。融合后的表示提升了路径排序和出行时间估计等任务的精度，包括少量或无标注数据的场景。}

\pubentry{AAAI, 2025, CCF A 类会议}{\href{https://ojs.aaai.org/index.php/AAAI/article/view/33350}{论文} \,/\, \href{https://arxiv.org/abs/2408.12809}{预印本}}
{DutyTTE: Deciphering Uncertainty in Origin-Destination Travel Time Estimation}
{Xiaowei Mao*, \me*, Shengnan Guo, Yubin Chen, Xingyu Xian, Haomin Wen, Qisen Xu, Youfang Lin, Huaiyu Wan}
{估计起点与终点之间的出行时间并量化每次预测的不确定性。该方法先推断可能采用的路线，再生成可靠的置信区间，以反映不同条件下各路段对整体不确定性的影响。}

\pubentry{SIGMOD, 2024, CCF A 类会议}{\href{https://dl.acm.org/doi/10.1145/3617337}{论文} \,/\, \href{https://github.com/Logan-Lin/DOT}{代码}}
{Origin-Destination Travel Time Oracle for Map-based Services}
{\me, Huaiyu Wan, Jilin Hu, Shengnan Guo, Bin Yang, Christian S. Jensen, Youfang Lin}
{从连接相同起终点的大量历史行程中学习，估计给定出发时间的起终点出行时间。该方法先推断两点之间的可能路线，再预测沿线出行时间，从而提升地图导航服务的估计精度。}

\subsection*{人工智能材料科学}
人工智能材料科学是我正在探索的新方向，其潜力在于用人工智能发现具有定制性能的新材料，克服传统材料设计方法面临的困难。

\pubentry{Advanced Materials, 2026, 影响因子 29.1}{\href{https://doi.org/10.1002/adma.202522493}{论文} \,/\, \href{https://arxiv.org/abs/2509.13916}{预印本} \,/\, \href{https://github.com/Logan-Lin/AMDEN-code}{代码}}
{Inverse Design of Amorphous Materials with Targeted Properties}
{Jonas A. Finkler*, \me*, Tao Du, Jilin Hu, Morten M. Smedskjaer}
{生成符合目标性能的玻璃等无序材料原子结构，并将其优化为稳定的低能构型。该工作还发布了新的非晶材料数据集，以支持此类设计研究。}

\pubentry{NeurIPS Workshop, 2025}{\href{https://openreview.net/forum?id=cEgjPFdLvl}{论文}}
{AMDEN: Amorphous Materials DEnoising Network}
{Jonas A. Finkler*, \me*, Morten M. Smedskjaer, Tao Du, Jilin Hu}{}

\pubentry{PNCS17, 口头报告, 2025}{\href{https://github.com/Logan-Lin/amden-presentation/releases/latest/download/slide.pdf}{幻灯片}}
{AMDEN: Amorphous Materials DEnoising Network}
{\me*, Jonas A. Finkler*, Tao Du, Morten M. Smedskjaer, Jilin Hu}{}

\pubentry{AAAI, 2026, CCF A 类会议}{\href{https://ojs.aaai.org/index.php/AAAI/article/view/37139}{论文} \,/\, \href{https://arxiv.org/abs/2511.12489}{预印本}}
{SculptDrug: A Spatial Condition-Aware Bayesian Flow Model for Structure-based Drug Design}
{Qingsong Zhong, Haomin Yu, \me, Wangmeng Shen, Long Zeng, Jilin Hu}
{生成适配目标蛋白质三维结构的药物分子，使分子契合目标蛋白质表面，并同时适配其整体形态和局部细节，从而为基于结构的药物发现生成更准确的候选分子。}

\pubentry{Villum Foundation, 项目}{}
{基于扩散概率模型的材料逆向设计}
{}
{用扩散概率模型理解材料的化学成分、结构和性能之间的关系，并逆向设计新材料。}
